\section{Introduction}
\label{sec:intro}

% source: PAPER_SKELETON Sec 1; reframed per README Sec 0 (E1 non-replication) and E1_ADJUDICATION.

\paragraph{A model that continues its own flawed reasoning can silently
absorb the flaw.} Language models generate long reasoning traces one token
at a time, and a false intermediate step, whether hallucinated, mis-copied,
or adversarially planted, does not announce itself. The question this paper
asks is behavioral and narrow: when a false statement already sits in a
model's own trace and the model keeps going with no external prompt to
check, does it reject the statement or reuse it. The answer depends on the
model. On Qwen2.5-7B-Instruct, the response is governed by two things,
whether the contradiction is directly visible in the trace and whether the
false step is useful for finishing the proof, and not by how ``checkable''
the error is in principle. Across four further open-weights models this
structure largely does not transfer: the visibility pattern reverses on
three of them, and absorption itself disappears at 32B scale. Two
regularities hold in every model we tested: none re-verifies a value that
must be recomputed, and verbal rejection is a poor guide to genuine
checking.

\paragraph{Evidence distance does not predict the response.} A natural
hypothesis is that rejection is graded by inferential distance to the
refuting evidence. A matched-distance test on Qwen2.5-7B-Instruct, laid out
as six distance-linked predictions, gives no support for it: five fail, and
the point estimates for the two that carry the gradient story run opposite to
it.
% source: EVIDENCE_protocol_prereg P5 (commit/sha/timestamp); PAPER_SKELETON 1.2
The one prediction the test supports is not a distance gradient but a
polarity control, that affirmative and negated false statements are absorbed
equivalently. What predicts the response instead is whether the contradiction
is directly visible and whether the false step is useful for finishing the
proof. The mechanics of how the predictions were fixed in advance are
collected in Appendix~\ref{app:provenance}.

\paragraph{Visibility and usability boundaries govern the anchor model.} On
the anchor model, explicit rejection of a planted falsehood
is abundant when the contradiction is a directly readable premise, roughly
0.25 of continuations at measured distance $d{=}0$, and collapses about
eightfold to 0.03 one hop away, staying flat for larger distances.
% source: EVIDENCE_findings_core F-vis.1 (0.2467 d0 vs 0.0311 d1); PAPER_SKELETON 4.1
Separately, silent downstream reuse of a planted falsehood tracks whether
that falsehood is derivationally needed to finish the current proof. A
usable falsehood is reused near 0.8 whether its refutation is one hop or
infinitely many hops away, and an inert one is reused near zero.
% source: EVIDENCE_findings_core F-use.1 (0.863/0.768/0.787 usable; ~0 inert)
Both are boundary conditions on when in-stream self-checking happens at
all, and neither is a graded distance gradient.

\paragraph{The boundary conditions do not replicate across models.} The
single-model-family base is the obvious weakness, so we ran a multi-model
replication that repeats the three boundary conditions on four more
open-weights models: Llama-3.1-8B, OLMo-2-7B, Qwen2.5-1.5B, and
Qwen2.5-32B, with its predictions again fixed before the data were
generated and its verdict independently re-checked from the raw data.
% source: E1_ADJUDICATION (registration bdb098a1); E1_VERIFICATION (CONFIRMED)
Zero of the four fully replicate. The visibility cliff is the sharpest
casualty: on three of the four models the effect significantly reverses, so
explicit rejection is \emph{lower} when the contradiction is directly
visible than when it is a hop away, the opposite of the anchor.
% source: E1_ADJUDICATION Sec 3 B1 (llama8b -0.144, qwen1.5b -0.142, qwen32b -0.069)
The usability asymmetry is the most portable of the three, holding on two
of the four, but it is scale-bounded: usable-error absorption falls from
0.71 at 1.5B to 0.81 at 7B to 0.01 at 32B, so the effect essentially
vanishes with capability within a single model family.
% source: E1_ADJUDICATION Sec 4.2 (0.711/0.807/0.013); F-use.1 E1 caveat
Because the same pipeline reproduces the Qwen2.5-7B anchor exactly, this is
a property of the models, not of the analysis. We therefore present the
single-model-family limitation not as an open caveat but as a positive
finding: how a model handles a planted falsehood in its own trace is
model-dependent and scale-dependent.

\paragraph{Two results do survive across scale and lineage.} First, a
recompute boundary. In a program-execution regime, no model we tested,
including two frontier models that reject re-readable planted errors near
totally, re-verifies a value it must recompute; recomputable falsehoods are
absorbed at 0.98 to 1.00 everywhere.
% source: EVIDENCE_lineage_boundaries L-rec.1 (deep_kc5 1.000 haiku/sonnet/qwen)
Second, an instrument caution that holds wherever we can check it. Verbal
``rejection'' cannot by itself distinguish a genuine check from lexical
cueing: a frequency-matched but non-refuting rule elicits the same
rejection rate as genuine refuting evidence, and every one of those
frequency-matched rejections fabricates a rule that is absent from the
world.
% source: EVIDENCE_findings_core F-lex.1/F-lex.2 (FREQ 0.118 ~ REFUTING 0.114; 26/26 fabricated)
Only a derivation-licensed split, which asks whether the model's own text
contains a valid derivation, separates the two.

\paragraph{Scope guard.} All of these findings concern non-deliberative,
in-stream continuation on synthetic first-order-logic worlds. We do not
claim they extend to explicit deliberation or thinking modes, to natural
prose, or to tasks outside the audited-logic substrate; models with a
separate deliberation channel are deliberately excluded from the
controlled tests.
% source: PAPER_SKELETON 1.4; EVIDENCE_lineage_boundaries L-prod.1 (instruct exclusion)
We also frame the in-stream-versus-fresh-turn part of the production
boundary as corroboration of concurrent work on the self-correction blind
spot, discussed in Section~\ref{sec:related}, not as a discovery of ours.

\paragraph{Contributions.} We contribute (i) an audited perturb-and-validate
instrument that truth-audits every plant, formally validates every
continuation, and measures refutation distance mechanically rather than by
design label (Section~\ref{sec:protocol}); (ii) evidence that inferential
distance does not predict in-stream rejection or reuse
(Section~\ref{sec:reversal}); (iii) two
boundary conditions, visibility and usability, on the anchor model
(Sections~\ref{sec:visibility}, \ref{sec:usability}), together with a
multi-model result that they largely do not replicate; (iv)
a recompute boundary that does survive across scale and lineage
(Section~\ref{sec:production}); and (v) an instrument caution that verbal
rejection is lexically cued (Section~\ref{sec:instrument}). We read the
whole as a boundary-condition map with an honest account of how far each
boundary travels.
